%% ============================================================
%%  Aula 4 — Controle sem modelo e aproximação de função valor
%%  Adaptação em pt-BR do CS234 (Stanford, Emma Brunskill), Lecture 4
%%  Compilar com XeLaTeX (2x) — latexmk -xelatex aula04.tex
%% ============================================================
\documentclass[aspectratio=169, 11pt]{beamer}
%% .sty do projeto na raiz do repositório (../)
\makeatletter
\def\input@path{{./}{../}}
\makeatother


\usetheme{AprendizadoRL}      % ANTES do babel — fontspec precisa vir primeiro
\usepackage{babel}
\babelprovide[import, main]{brazilian}
\usepackage{booktabs}
\usepackage{multirow}

\title[Aula 4 — Controle sem modelo]{Aula 4: Controle sem modelo\\ e aproximação de função valor}
\subtitle[Aprendizagem por Reforço]{MC/TD em lote \textbullet\ $\epsilon$-guloso e GLIE \textbullet\ SARSA e Q-learning \textbullet\ VFA \textbullet\ DQN}
\author[Prof. Marcelo K. Albertini]{Prof. Marcelo Keese Albertini \\
  \footnotesize Universidade Federal de Uberlândia \\
  \footnotesize adaptado de Emma Brunskill, CS234 (Stanford)}
\institute{Universidade Federal de Uberlândia}
\date{2026}

% ---- Nota discreta: perguntas ficam para o professor conduzir em aula
\newcommand{\paradiscussao}{\smallskip
  {\footnotesize\color{rlCinza}\itshape Pergunta para discussão conduzida em aula.}}

% ---- Notação local desta aula
\newcommand{\Qhat}{\widehat{Q}}
\newcommand{\Vhat}{\widehat{V}}
\newcommand{\Phat}{\widehat{P}}
\newcommand{\rhat}{\widehat{\rec}}
\newcommand{\wv}{\mathbf{w}}
\newcommand{\wtar}{\mathbf{w}^{-}}
\newcommand{\xv}{\mathbf{x}}
\newcommand{\polb}{\pol_b}                 % política de comportamento
\newcommand{\Qstar}{Q^{*}}

% ---- Pseudocódigo compacto e numerado
\newcounter{plin}
\newcommand{\pl}[2][0em]{%
  \stepcounter{plin}%
  \par\noindent
  \hangindent=2.0em \hangafter=1
  \makebox[1.5em][r]{\color{rlCinza}\scriptsize\theplin:}\hspace{0.4em}\hspace*{#1}\ignorespaces#2\par}
\newenvironment{pseudo}%
  {\par\setcounter{plin}{0}\footnotesize\setlength{\parskip}{0.25ex}}%
  {\par}

\begin{document}

\begin{frame}[plain, noframenumbering]\titlepage\end{frame}

%% ------------------------------------------------------------
\begin{frame}{Plano de hoje}
  \framesubtitle{A aula em que o agente deixa de observar e passa a decidir}

  \begin{enumerate}\setlength{\itemsep}{0.55ex}
    \item \textbf{Avaliação de política em lote} — fechamos a Aula 3
      entendendo \emph{o que} MC e TD calculam quando os dados acabam.
    \item \textbf{Melhoria de política generalizada} — por que precisamos
      de $Q$ e não de $V$, e como explorar sem perder a garantia de
      melhoria monotônica.
    \item \textbf{Controle tabular sem modelo} — MC com $\epsilon$-guloso,
      GLIE, SARSA e Q-learning.
    \item \textbf{Aproximação de função valor} — abandonar a tabela para
      \emph{generalizar} entre estados.
    \item \textbf{Deep Q-Learning} — o que foi preciso acrescentar para
      que a ideia funcionasse em Atari.
  \end{enumerate}

  \begin{ideiabox}
    Fio condutor: os quatro algoritmos de hoje são a mesma linha de código,
    $\;$\emph{estimativa} $\leftarrow$ \emph{estimativa} $+\;\alpha\,($\emph{alvo} $-$ \emph{estimativa}$)$,
    com quatro escolhas diferentes de \emph{alvo}.
  \end{ideiabox}
\end{frame}

%% ------------------------------------------------------------
\begin{frame}{Onde estamos}
  \framesubtitle{Duas dimensões: temos o modelo? queremos avaliar ou decidir?}

  \begin{center}\small
  \begin{tabular}{lll}
    \toprule
     & \textbf{Avaliar $\pol$} (predição) & \textbf{Escolher $\pol$} (controle) \\
    \midrule
    \textbf{Com modelo} & Avaliação de política (Aula 2) & Iteração de política / de valor (Aula 2) \\
    \addlinespace[0.3ex]
    \textbf{Sem modelo} & MC, TD(0) (Aula 3) & \alert{MC-controle, SARSA, Q-learning (hoje)} \\
    \bottomrule
  \end{tabular}
  \end{center}

  \medskip
  \begin{itemize}\setlength{\itemsep}{0.3ex}
    \item Até agora tudo era \textbf{tabular}: uma entrada de memória por
      estado (ou por par estado-ação).
    \item A segunda metade da aula quebra essa hipótese: $\abs{\gS}$ grande
      demais para caber em memória, quanto mais para ser visitado.
  \end{itemize}
\end{frame}

%% ------------------------------------------------------------
\begin{frame}{Recapitulação: os três estimadores da Aula 3}
  \framesubtitle{Estimar $\Vpi(s)$ executando $\pol$, sem conhecer $P$ e $\rec$}

  \begin{itemize}\setlength{\itemsep}{0.45ex}
    \item \textbf{Monte Carlo} — precisa da trajetória completa $\tau$ e usa
      o retorno observado $G_t$:
      \[ \Vpi(s) \leftarrow \Vpi(s) + \alpha\bigl(G_t(s) - \Vpi(s)\bigr) \]
    \item \textbf{TD(0)} — precisa apenas da tupla $(s,a,r,s')$ e
      \emph{bootstrapa} na própria estimativa:
      \[ \Vpi(s) \leftarrow \Vpi(s) + \alpha\bigl(\underbrace{r + \gamma \Vpi(s')}_{\text{alvo TD}} - \Vpi(s)\bigr) \]
    \item \textbf{Equivalência de certeza} — usa as tuplas para estimar
      $\Phat$ e $\rhat$ por máxima verossimilhança e roda programação
      dinâmica sobre esse MDP estimado.
  \end{itemize}

  \begin{ideiabox}
    Os três já apareceram como \emph{procedimentos}. Hoje começamos
    perguntando algo mais fino: para \emph{qual objeto matemático} cada um
    converge quando os dados são finitos e reciclados?
  \end{ideiabox}
\end{frame}

%% ============================================================
\section{Avaliação de política em lote}
%% ============================================================

\begin{frame}{Do fluxo para o lote}
  \framesubtitle{\emph{Batch} (ou \emph{offline}): o conjunto de dados é fixo}

  \begin{itemize}\setlength{\itemsep}{0.4ex}
    \item Nas versões vistas até aqui, cada tupla $(s,a,r,s')$ é usada
      \textbf{uma vez} e descartada.
    \item No cenário em lote temos $K$ episódios armazenados e podemos
      reprocessá-los quantas vezes quisermos.
  \end{itemize}

  \begin{algbox}{Avaliação de política em lote}
    \begin{pseudo}
      \pl{Dado o conjunto fixo de $K$ episódios}
      \pl{\textbf{repita} até convergir}
      \pl[1.2em]{amostre um episódio $k \sim \{1,\dots,K\}$}
      \pl[1.2em]{aplique a atualização de MC (ou de TD(0)) a esse episódio}
      \pl{\textbf{fim}}
    \end{pseudo}
  \end{algbox}

\end{frame}

%% ------------------------------------------------------------
\begin{frame}{O exemplo A--B}
  \framesubtitle{Sutton \& Barto (2018), Exemplo 6.4 — dois estados, $\gamma = 1$}

  \begin{columns}[T]
    \begin{column}{0.44\textwidth}
      \textbf{Oito episódios observados:}
      \smallskip

      {\small
      \begin{tabular}{ll}
        \toprule
        Episódio & Frequência \\
        \midrule
        $A,\,0,\,B,\,0$ & 1 vez \\
        $B,\,1$         & 6 vezes \\
        $B,\,0$         & 1 vez \\
        \bottomrule
      \end{tabular}}

      \medskip
      {\small Cada episódio termina após a última recompensa. Note que
      $A$ aparece \textbf{uma única vez} em toda a base.}
    \end{column}
    \begin{column}{0.54\textwidth}
      \begin{defbox}{}
        Estimativa a partir de $B$: em $8$ visitas houve recompensa $1$ em
        $6$ delas, logo $V(B) = 6/8 = 0{,}75$ tanto por MC quanto por TD.
      \end{defbox}

      \begin{quizbox}{}
        E $V(A)$? Rodando as atualizações sobre esses dados um número
        infinito de vezes, para que valor cada método converge?
      \end{quizbox}
      \paradiscussao
    \end{column}
  \end{columns}
\end{frame}

%% ------------------------------------------------------------
\begin{frame}{O exemplo A--B: as duas respostas}
  \framesubtitle{Não é um erro de um dos métodos — eles otimizam critérios diferentes}

  \begin{columns}[T]
    \begin{column}{0.5\textwidth}
      \begin{respbox}
        \[ V_{\mathrm{MC}}(A) = 0 \qquad V_{\mathrm{TD}}(A) = 0{,}75 \]
      \end{respbox}

      \smallskip
      \textbf{Monte Carlo.} O único retorno já observado a partir de $A$
      foi $G = 0 + 0 = 0$. MC minimiza o erro quadrático em relação aos
      retornos \emph{observados}, então responde $0$.

      \smallskip
      \textbf{TD(0).} O alvo é $r + \gamma V(B) = 0{,}75$: TD encadeia a
      informação obtida em $B$, mesmo que $A$ apareça uma vez só.
    \end{column}
    \begin{column}{0.5\textwidth}
      \begin{center}
      \scalebox{0.92}{\begin{tikzpicture}[node distance=16mm]
        \node[estado destacado] (a) {$A$};
        \node[estado, right=of a] (b) {$B$};
        \node[estado terminal, right=of b] (t) {$\bot$};
        \draw[trans] (a) -- node[probl] {$1{,}0$} node[above, font=\scriptsize, text=rlCinza] {$r=0$} (b);
        \draw[trans] (b) -- node[probl] {$1{,}0$} node[above, font=\scriptsize, text=rlVerde!70!black] {$\E[r]=0{,}75$} (t);
      \end{tikzpicture}}
      \end{center}

      \vspace{-0.7em}
      \begin{ideiabox}
        Esse é o MDP de \textbf{máxima verossimilhança} implícito nos dados.
        TD(0) em lote converge exatamente para o $\Vpi$ desse MDP: por isso
        $V(A) = 0 + 0{,}75$.
      \end{ideiabox}
    \end{column}
  \end{columns}
\end{frame}

%% ------------------------------------------------------------
\begin{frame}{Para onde MC e TD convergem no cenário em lote}

  \begin{teobox}{Monte Carlo em lote}
    Converge para os valores que \textbf{minimizam o erro quadrático médio}
    em relação aos retornos observados,
    $V = \argmin_{V} \sum_{k,t} \bigl(G_{k,t} - V(s_{k,t})\bigr)^2$.
    No exemplo A--B: $V(A) = 0$.
  \end{teobox}

  \begin{teobox}{TD(0) em lote}
    Converge para o $\Vpi$ do MDP de máxima verossimilhança estimado dos
    dados --- ou seja, é \textbf{idêntico à equivalência de certeza}:
    {\small
    \[
      \Phat(s' \mid s,a) = \frac{\sum_{k} \mathbf{1}\{s_k=s,\,a_k=a,\,s_{k+1}=s'\}}{N(s,a)},
      \qquad
      \rhat(s,a) = \frac{\sum_{k} \mathbf{1}\{s_k=s,\,a_k=a\}\,r_k}{N(s,a)}
    \]}
    \vspace{-0.7em}
    No exemplo A--B: $V(A) = 0{,}75$.
  \end{teobox}
\end{frame}

%% ------------------------------------------------------------
\begin{frame}{Qual das duas respostas é ``a certa''?}
  \framesubtitle{Depende inteiramente de acreditarmos, ou não, na hipótese de Markov}

  \begin{columns}[T]
    \begin{column}{0.5\textwidth}
      \begin{ideiabox}
        \textbf{Se o ambiente é markoviano}, TD está certo: o que se aprende
        sobre $B$ vale para \emph{qualquer} caminho que leve a $B$. TD usa a
        estrutura; MC a ignora.
      \end{ideiabox}
    \end{column}
    \begin{column}{0.5\textwidth}
      \begin{alertabox}
        \textbf{Se o ambiente não é markoviano} (estado parcialmente
        observado, por exemplo), TD propaga uma hipótese falsa. MC, que só
        olha retornos completos, é mais robusto — ao custo de variância.
      \end{alertabox}
    \end{column}
  \end{columns}

  \smallskip
  {\small MC é correto \emph{sem} supor Markov; TD supõe Markov e, quando
  ela vale, converge com muito menos dados.}

  \vspace{-0.4em}
  \begin{quizbox}{}
    Com prontuários médicos incompletos, você confiaria mais em MC ou em
    TD? Por quê?
  \end{quizbox}

\end{frame}

%% ------------------------------------------------------------
\begin{frame}{Critérios para comparar algoritmos de avaliação}

  \begin{columns}[T]
    \begin{column}{0.52\textwidth}
      \textbf{Custo computacional}
      \begin{itemize}\setlength{\itemsep}{0.2ex}
        \item TD(0): $O(1)$ por transição, $O(L)$ por episódio.
        \item MC: também $O(L)$, mas só \emph{depois} do fim do episódio.
        \item Equivalência de certeza: resolve um sistema linear a cada
          atualização --- muito mais caro.
      \end{itemize}

      \smallskip
      \textbf{Eficiência de dados}
      \begin{itemize}\setlength{\itemsep}{0.2ex}
        \item Usar a estrutura de Markov (TD, eq.\ de certeza) extrai mais
          de cada amostra; MC pode ser preferível quando ela não vale.
      \end{itemize}
    \end{column}
    \begin{column}{0.48\textwidth}
      \begin{center}\footnotesize
      \begin{tabular}{lccc}
        \toprule
        & MC & TD(0) & Eq.\ cert. \\
        \midrule
        Usa Markov        & não & sim & sim \\
        Viés              & não & sim & sim \\
        Variância         & alta & baixa & baixa \\
        Episódio completo & sim & não & não \\
        Custo/atualiz.    & $O(1)$ & $O(1)$ & alto \\
        Eficiência dados  & baixa & média & alta \\
        \bottomrule
      \end{tabular}
      \end{center}

    \end{column}
  \end{columns}
\end{frame}

%% ------------------------------------------------------------
\begin{frame}{Fecho da avaliação de política}
  \framesubtitle{Estimar o retorno esperado de uma $\pol$ dada, sem o MDP verdadeiro}

  \begin{itemize}\setlength{\itemsep}{0.4ex}
    \item \textbf{Monte Carlo} — com dados \emph{on-policy} ou
      \emph{off-policy} (amostragem por importância).
    \item \textbf{Diferença temporal} — TD(0) e suas variantes.
    \item \textbf{Programação dinâmica com equivalência de certeza} —
      construir o modelo e resolvê-lo.
  \end{itemize}

  \vspace{-0.4em}
  \begin{defbox}{A régua que usaremos o curso inteiro}
    Supõe Markov? \; É consistente? \; Qual o viés e a variância? \;
    Quantos dados precisa? \; Quanto custa por atualização?
  \end{defbox}

  \begin{ideiabox}
    Avaliação é só metade do problema. A partir de agora queremos
    \emph{escolher} a política, não apenas medi-la --- e aí aparece um
    obstáculo que não existia até aqui: precisamos de dados sobre ações
    que a política atual \textbf{não} toma.
  \end{ideiabox}
\end{frame}

%% ============================================================
\section{Melhoria de política generalizada}
%% ============================================================

\begin{frame}{Iteração de política sem modelo}
  \framesubtitle{O mesmo esqueleto da Aula 2 --- com três buracos}

  \begin{columns}[T]
    \begin{column}{0.54\textwidth}
      \vspace{-0.3em}
      \begin{algbox}{Iteração de política generalizada}
        \begin{pseudo}
          \pl{Inicialize $\pol$}
          \pl{\textbf{repita}}
          \pl[1.2em]{\emph{Avaliação:} estime $\Qpi$}
          \pl[1.2em]{\emph{Melhoria:} atualize $\pol$ a partir de $\Qpi$}
          \pl{\textbf{fim}}
        \end{pseudo}
      \end{algbox}

      \vspace{-0.3em}
      \begin{alertabox}
        \begin{enumerate}\setlength{\itemsep}{0.15ex}
          \item Se $\pol$ é determinística, $Q(s,a)$ com $a \neq \pol(s)$
            nunca é observado.
          \item A melhoria usa um $\Qhat$ estimado, com ruído.
          \item Quanto avaliar antes de melhorar?
        \end{enumerate}
      \end{alertabox}
    \end{column}
    \begin{column}{0.46\textwidth}
      \begin{center}
      \scalebox{0.74}{\begin{tikzpicture}[node distance=8mm]
        \node[draw=rlAzul, very thick, fill=rlAzul!8, rounded corners=2pt,
              minimum width=32mm, minimum height=9mm, font=\small] (av) {avaliação: $\Qhat \approx \Qpi$};
        \node[draw=rlLaranja, very thick, fill=rlLaranja!12, rounded corners=2pt,
              minimum width=32mm, minimum height=9mm, font=\small, below=14mm of av] (me)
              {melhoria: $\pol \leftarrow$ guloso$(\Qhat)$};
        \draw[trans destac] ([xshift=-13mm]av.south) to[bend right=38] ([xshift=-13mm]me.north);
        \draw[trans destac] ([xshift=13mm]me.north) to[bend right=38] ([xshift=13mm]av.south);
        \node[font=\scriptsize, text=rlCinza, right=0mm of me.south east, anchor=north east,
              yshift=-3mm, text width=36mm, align=center] {converge para $\polstar$ e $\Qstar$};
      \end{tikzpicture}}
      \end{center}
    \end{column}
  \end{columns}
\end{frame}

%% ------------------------------------------------------------
\begin{frame}{Por que $Q$ e não $V$}
  \framesubtitle{Uma consequência silenciosa de não ter o modelo}

  \begin{itemize}\setlength{\itemsep}{0.4ex}
    \item Com modelo, a melhoria gulosa se faz a partir de $V$:
      \[ \pol'(s) = \argmax_{a \in \gA}\; \Bigl[\rec(s,a) + \gamma \sum_{s'} \tran{s'}{s,a} V(s')\Bigr] \]
    \item Sem modelo, $\rec$ e $P$ são desconhecidos --- a expressão acima é
      inútil. Já com $Q$:
      \[ \pol'(s) = \argmax_{a \in \gA}\; Q(s,a) \]
      e a melhoria é imediata, sem nenhum conhecimento da dinâmica.
  \end{itemize}

  \begin{ideiabox}
    $Q$ é exatamente $V$ com a \emph{informação de um passo do modelo já
    embutida}. É por isso que praticamente todo algoritmo de controle sem
    modelo aprende $Q$, e não $V$.
  \end{ideiabox}

\end{frame}

%% ------------------------------------------------------------
\begin{frame}{O problema da exploração}
  \framesubtitle{O dilema que separa RL de aprendizado supervisionado}

  \begin{columns}[T]
    \begin{column}{0.5\textwidth}
      \textbf{Explorar}
      \begin{itemize}\setlength{\itemsep}{0.2ex}
        \item Não há como aprender sobre uma ação sem experimentá-la.
        \item Dados sobre $a$ só existem se $\pol$ der probabilidade
          positiva a $a$.
      \end{itemize}
    \end{column}
    \begin{column}{0.5\textwidth}
      \textbf{Explotar}
      \begin{itemize}\setlength{\itemsep}{0.2ex}
        \item Cada passo experimentando é um passo não gasto colhendo
          recompensa conhecida.
        \item Em saúde e educação, explorar tem custo real.
      \end{itemize}
    \end{column}
  \end{columns}

  \vspace{-0.2em}
  {\small Uma política \textbf{gulosa pura} sobre $\Qhat$ pode travar para
  sempre: basta que a primeira amostra de uma ação boa seja ruim por azar
  para que ela nunca mais seja escolhida --- nem corrigida.}

  \vspace{-0.4em}
  \begin{quizbox}{}
    O rover tenta ``ir para a direita'', recebe $0$ por azar. O que faz
    uma política gulosa pura daí em diante?
  \end{quizbox}
\end{frame}

\begin{frame}{Políticas $\epsilon$-gulosas}
  \framesubtitle{A solução mais simples que funciona --- e a mais usada na prática}

  \begin{defbox}{Política $\epsilon$-gulosa em relação a $Q(s,a)$}
    Seja $\abs{\gA}$ o número de ações e $\epsilon \in [0,1]$. Então
    \[
      \pol(a \mid s) =
      \begin{cases}
        1 - \epsilon + \dfrac{\epsilon}{\abs{\gA}}, & \text{se } a = \argmax_{a'} Q(s,a')\\[1.6ex]
        \dfrac{\epsilon}{\abs{\gA}}, & \text{caso contrário.}
      \end{cases}
    \]
  \end{defbox}

  \begin{itemize}\setlength{\itemsep}{0.3ex}
    \item Toda ação tem probabilidade $\ge \epsilon/\abs{\gA} > 0$ ---
      exploração garantida.
    \item $\epsilon = 1$: passeio aleatório. \quad $\epsilon = 0$: guloso puro.
    \item Alternativa comum: \emph{softmax}, $\pol(a\mid s) \propto e^{Q(s,a)/\tau}$,
      que explora preferencialmente as ações \emph{quase} boas.
  \end{itemize}
\end{frame}

%% ------------------------------------------------------------
\begin{frame}{A melhoria monotônica sobrevive à exploração}

  \begin{itemize}\setlength{\itemsep}{0.35ex}
    \item Na Aula 2 provamos que a iteração de política, com $P$ e $\rec$
      conhecidos, melhora monotonicamente --- mas a prova supunha que a
      melhoria produzisse uma política \textbf{determinística}.
    \item Uma política $\epsilon$-gulosa é estocástica de propósito. A
      garantia se perde?
  \end{itemize}

  \begin{teobox}{Teorema — melhoria monotônica $\epsilon$-gulosa}
    Seja $\pol_i$ uma política $\epsilon$-gulosa qualquer e seja
    $\pol_{i+1}$ a política $\epsilon$-gulosa em relação a $Q^{\pol_i}$.
    Então
    \[ V^{\pol_{i+1}}(s) \;\ge\; V^{\pol_i}(s) \qquad \forall s \in \gS. \]
  \end{teobox}

  \vspace{-0.4em}
  \begin{ideiabox}
    Ou seja: é possível explorar \emph{para sempre} e ainda assim nunca
    piorar. O custo é que a política ótima entre as $\epsilon$-gulosas não
    é $\polstar$ --- ela carrega o $\epsilon$ de ruído. É esse resíduo que a
    condição GLIE, adiante, elimina.
  \end{ideiabox}
\end{frame}

%% ------------------------------------------------------------
\begin{frame}{Prova da melhoria monotônica $\epsilon$-gulosa}
  \framesubtitle{O truque é escrever o máximo como uma média ponderada disfarçada}

  {\footnotesize
  \begin{align*}
    Q^{\pol_i}\bigl(s, \pol_{i+1}(s)\bigr)
      &= \sum_{a \in \gA} \pol_{i+1}(a\mid s)\, Q^{\pol_i}(s,a) \\
      &= \frac{\epsilon}{\abs{\gA}} \sum_{a \in \gA} Q^{\pol_i}(s,a)
         \;+\; (1-\epsilon)\, \max_{a} Q^{\pol_i}(s,a) \\[0.4ex]
      &\ge \frac{\epsilon}{\abs{\gA}} \sum_{a \in \gA} Q^{\pol_i}(s,a)
         \;+\; (1-\epsilon) \sum_{a \in \gA}
           \underbrace{\frac{\pol_i(a\mid s) - \epsilon/\abs{\gA}}{1-\epsilon}}_{\text{pesos} \;\ge 0,\; \text{somam } 1}
           Q^{\pol_i}(s,a) \\[0.4ex]
      &= \sum_{a \in \gA} \pol_i(a \mid s)\, Q^{\pol_i}(s,a) \;=\; V^{\pol_i}(s).
  \end{align*}}

  \vspace{-0.5em}
  \begin{itemize}\setlength{\itemsep}{0.2ex}
    \item[] {\footnotesize A desigualdade vale porque o máximo é maior ou
      igual a \emph{qualquer} combinação convexa. Os pesos são não-negativos
      precisamente porque $\pol_i$ é $\epsilon$-gulosa, logo
      $\pol_i(a\mid s) \ge \epsilon/\abs{\gA}$.}
    \item[] {\footnotesize De $Q^{\pol_i}(s,\pol_{i+1}(s)) \ge V^{\pol_i}(s)$
      segue $V^{\pol_{i+1}} \ge V^{\pol_i}$ pelo teorema de melhoria de
      política da Aula 2. \hfill $\blacksquare$}
  \end{itemize}
\end{frame}

%% ============================================================
\section{Controle tabular sem modelo}
%% ============================================================

\begin{frame}{Monte Carlo para $Q$}
  \framesubtitle{A avaliação da Aula 3, com o par $(s,a)$ no lugar de $s$}

  \begin{algbox}{Avaliação de política Monte Carlo para $Q^{\pol}$ (primeira visita)}
    \begin{pseudo}
      \pl{Inicialize $Q(s,a) = 0$, $N(s,a)=0$ para todo $(s,a)$; $k = 1$; política $\pol$}
      \pl{\textbf{repita}}
      \pl[1.2em]{gere o $k$-ésimo episódio $(s_{k,1}, a_{k,1}, r_{k,1}, \dots, s_{k,T})$ seguindo $\pol$}
      \pl[1.2em]{calcule $G_{k,t} = r_{k,t} + \gamma r_{k,t+1} + \gamma^2 r_{k,t+2} + \cdots$ para todo $t$}
      \pl[1.2em]{\textbf{para} $t = 1,\dots,T$ \textbf{faça}}
      \pl[2.4em]{\textbf{se} é a primeira visita a $(s_t,a_t)$ no episódio $k$ \textbf{então}}
      \pl[3.6em]{$N(s_t,a_t) \mathrel{+}= 1$}
      \pl[3.6em]{$Q(s_t,a_t) \mathrel{+}= \frac{1}{N(s_t,a_t)}\bigl(G_{k,t} - Q(s_t,a_t)\bigr)$}
      \pl[2.4em]{\textbf{fim}}
      \pl[1.2em]{\textbf{fim}; \quad $k \mathrel{+}= 1$}
      \pl{\textbf{fim}}
    \end{pseudo}
  \end{algbox}

\end{frame}

%% ------------------------------------------------------------
\begin{frame}{Controle Monte Carlo \emph{on-policy}}
  \framesubtitle{Avaliar e melhorar a cada episódio, com $\epsilon$ decrescente}

  \begin{algbox}{Controle MC on-policy com $\epsilon$-guloso}
    \begin{pseudo}
      \pl{Inicialize $Q(s,a)=0$, $N(s,a)=0$ para todo $(s,a)$; $\epsilon = 1$; $k = 1$}
      \pl{$\pol_k \leftarrow \epsilon\text{-guloso}(Q)$ \hfill \textit{// política inicial}}
      \pl{\textbf{repita}}
      \pl[1.2em]{gere o episódio $k$ seguindo $\pol_k$ e calcule os retornos $G_{k,t}$}
      \pl[1.2em]{\textbf{para} $t = 1,\dots,T$ \textbf{faça}}
      \pl[2.4em]{\textbf{se} primeira visita a $(s_t,a_t)$ \textbf{então} $N \mathrel{+}= 1$;\;
                 $Q(s_t,a_t) \mathrel{+}= \frac{1}{N(s_t,a_t)}(G_{k,t} - Q(s_t,a_t))$}
      \pl[1.2em]{\textbf{fim}}
      \pl[1.2em]{$k \mathrel{+}= 1$; \quad $\epsilon \leftarrow 1/k$ \hfill \textit{// exploração decrescente}}
      \pl[1.2em]{$\pol_k \leftarrow \epsilon\text{-guloso}(Q)$ \hfill \textit{// melhoria de política}}
      \pl{\textbf{fim}}
    \end{pseudo}
  \end{algbox}

\end{frame}

%% ------------------------------------------------------------
\begin{frame}{Exemplo trabalhado: MC-controle no \emph{Mars Rover}}
  \framesubtitle{Sete estados, duas ações, $\gamma = 1$}

  \begin{center}\scalebox{0.68}{\marsrover{1,7}}\end{center}

  \smallskip
  \begin{itemize}\setlength{\itemsep}{0.25ex}
    \item Recompensas: $\rec(\cdot, a_1) = [\,1\;0\;0\;0\;0\;0\;{+}10\,]$,\quad
      $\rec(\cdot, a_2) = [\,0\;0\;0\;0\;0\;0\;{+}5\,]$.
    \item Início: $Q(s,a) = 0$ para todo par, $\pol(s) = a_1$, $\epsilon = 0{,}5$.
    \item Trajetória amostrada da política $\epsilon$-gulosa:
      \[ \tau = (s_3, a_1, 0,\; s_2, a_2, 0,\; s_3, a_1, 0,\; s_2, a_2, 0,\; s_1, a_1, 1,\; \bot) \]
  \end{itemize}

  \begin{quizbox}{}
    Qual é a estimativa MC de primeira visita de $Q$ para cada par
    $(s,a)$ após essa única trajetória? E qual passa a ser a política
    gulosa?
  \end{quizbox}
  \paradiscussao
\end{frame}

%% ------------------------------------------------------------
\begin{frame}{Exemplo trabalhado: solução}
  \framesubtitle{Com $\gamma = 1$, todo retorno é a soma das recompensas até o fim}

  \begin{columns}[T]
    \begin{column}{0.52\textwidth}
      \begin{respbox}
        \begin{itemize}\setlength{\itemsep}{0.25ex}
          \item $Q(\cdot, a_1) = [\,1\;\;0\;\;1\;\;0\;\;0\;\;0\;\;0\,]$
          \item $Q(\cdot, a_2) = [\,0\;\;1\;\;0\;\;0\;\;0\;\;0\;\;0\,]$
          \item Nova política gulosa: $\pol = [\,a_1\;\;a_2\;\;a_1\;\;$empate$\;\cdots]$
        \end{itemize}
      \end{respbox}

      \smallskip
      {\footnotesize Toda primeira visita ocorre antes da recompensa $1$
      em $s_1$: com $\gamma=1$, cada par visitado recebe retorno $1$.}
    \end{column}
    \begin{column}{0.48\textwidth}
    
      \begin{quizbox}{}
        Com $\epsilon = 1/3$ e $\abs{\gA} = 2$, qual a probabilidade de a
        nova política escolher $a_1$ em $s_1$?
      \end{quizbox}

      \pause
      \begin{respbox}
        $\;1 - \epsilon + \dfrac{\epsilon}{\abs{\gA}}
          = 1 - \dfrac{1}{3} + \dfrac{1}{6} = \dfrac{5}{6}$.

        \smallskip
        {\footnotesize Erro comum: responder $2/3$, esquecendo que a
        exploração uniforme também pode cair na ação gulosa.}
      \end{respbox}
    \end{column}
  \end{columns}
\end{frame}

%% ------------------------------------------------------------
\begin{frame}{Guloso no limite com exploração infinita (GLIE)}
  \framesubtitle{A condição que transforma ``nunca piora'' em ``converge para o ótimo''}

  \begin{defbox}{GLIE (\emph{Greedy in the Limit with Infinite Exploration})}
    Uma sequência de políticas $\{\pol_i\}$ é GLIE se
    \begin{enumerate}\setlength{\itemsep}{0.3ex}
      \item todo par $(s,a)$ é visitado infinitas vezes,
        $\lim_{i \to \infty} N_i(s,a) \to \infty$;
      \item a política de comportamento converge para a gulosa,
        $\lim_{i \to \infty} \pol_i(a \mid s) \to \argmax_{a} Q(s,a)$ com
        probabilidade $1$.
    \end{enumerate}
  \end{defbox}

  {\small As duas condições puxam em direções opostas. A
  \textbf{estratégia GLIE mais simples} concilia ambas: $\epsilon$-guloso
  com $\epsilon_i = 1/i$, pois $\sum_i 1/i$ diverge mas $\epsilon_i \to 0$.}

  \begin{teobox}{Teorema — convergência do controle MC com GLIE}
    O controle Monte Carlo GLIE converge para a função valor de ação ótima:
    $Q(s,a) \to \Qstar(s,a)$.
  \end{teobox}
\end{frame}

%% ------------------------------------------------------------
\begin{frame}{Controle com métodos de diferença temporal}
  \framesubtitle{Mesma receita, mas agora aprendendo a cada \emph{passo}}

  \begin{algbox}{Iteração de política generalizada com TD}
    \begin{pseudo}
      \pl{Inicialize $\pol$ (e $Q$)}
      \pl{\textbf{repita}}
      \pl[1.2em]{\emph{Avaliação:} atualize $Q$ por diferença temporal usando dados da $\epsilon$-gulosa}
      \pl[1.2em]{\emph{Melhoria:} $\pol \leftarrow \epsilon\text{-guloso}(Q)$}
      \pl{\textbf{fim}}
    \end{pseudo}
  \end{algbox}

  {\small Vantagem imediata: não é preciso esperar o episódio terminar ---
  funciona em tarefas \textbf{contínuas}. E a pergunta que separa os dois
  algoritmos de hoje: \textbf{qual $Q$ do próximo estado entra no alvo?}}

  \vspace{-0.4em}
  \begin{ideiabox}
    \textbf{SARSA} usa o valor da ação que o agente \emph{de fato tomou}.
    \textbf{Q-learning} usa o valor da \emph{melhor} ação disponível,
    tenha ela sido tomada ou não. Uma diferença de um símbolo, com
    consequências grandes.
  \end{ideiabox}
\end{frame}

%% ------------------------------------------------------------
\begin{frame}{\emph{On-policy} e \emph{off-policy}}

  \begin{columns}[T]
    \begin{column}{0.5\textwidth}
      \begin{defbox}{Aprendizado on-policy}
        Estima e melhora a política \emph{a partir da experiência gerada
        por ela mesma}. A política que age é a política que se aprende.
        \smallskip

        {\footnotesize Exemplos: controle MC on-policy, SARSA.}
      \end{defbox}
    \end{column}
    \begin{column}{0.5\textwidth}
      \begin{defbox}{Aprendizado off-policy}
        Estima e melhora uma política \emph{alvo} usando dados coletados
        por uma política de \emph{comportamento} diferente, $\polb$.
        \smallskip

        {\footnotesize Exemplos: Q-learning, DQN, praticamente todo RL
        offline moderno.}
      \end{defbox}
    \end{column}
  \end{columns}

  \medskip
  \begin{itemize}\setlength{\itemsep}{0.3ex}
    \item Off-policy é o que permite aprender de \textbf{dados de arquivo}:
      registros médicos, logs de recomendação, demonstrações humanas.
    \item Também permite reutilizar experiência antiga --- a base do
      \emph{replay} do DQN --- mas é a fonte de boa parte da instabilidade
      que veremos ao combiná-lo com aproximação de função.
  \end{itemize}
\end{frame}

%% ------------------------------------------------------------
\begin{frame}{SARSA}
  \framesubtitle{O nome vem da quíntupla usada na atualização: $(s, a, r, s', a')$}

  \begin{defbox}{Atualização SARSA}
    \vspace{-1.2em}
    \[ Q(s_t, a_t) \leftarrow Q(s_t,a_t) + \alpha\Bigl(
       \underbrace{r_t + \gamma\, Q(s_{t+1}, a_{t+1})}_{\text{alvo: ação realmente tomada}} - Q(s_t,a_t)\Bigr) \]
  \end{defbox}

  \begin{algbox}{SARSA com exploração $\epsilon$-gulosa}
    \begin{pseudo}
      \pl{Inicialize $Q(s,a)$; $t=0$; estado inicial $s_0$; $\pol \leftarrow \epsilon\text{-guloso}(Q)$}
      \pl{Tome $a_t \sim \pol(s_t)$}
      \pl{\textbf{repita}}
      \pl[1.2em]{observe $(r_t, s_{t+1})$ e tome $a_{t+1} \sim \pol(s_{t+1})$}
      \pl[1.2em]{$Q(s_t,a_t) \mathrel{+}= \alpha\bigl(r_t + \gamma Q(s_{t+1},a_{t+1}) - Q(s_t,a_t)\bigr)$}
      \pl[1.2em]{$\pol \leftarrow \epsilon\text{-guloso}(Q)$; \quad $t \mathrel{+}= 1$}
      \pl{\textbf{fim}}
    \end{pseudo}
  \end{algbox}

\end{frame}

%% ------------------------------------------------------------
\begin{frame}{Q-learning}
  \framesubtitle{Aprender $\Qstar$ enquanto se age com outra política}

  \begin{defbox}{Atualização Q-learning \textnormal{(Watkins, 1989)}}
    \vspace{-1.2em}
    \[ Q(s_t, a_t) \leftarrow Q(s_t,a_t) + \alpha\Bigl(
       \underbrace{r_t + \gamma \max_{a'} Q(s_{t+1}, a')}_{\text{alvo: melhor ação possível}} - Q(s_t,a_t)\Bigr) \]
  \end{defbox}

  \vspace{-0.6em}
  \begin{algbox}{Q-learning com exploração $\epsilon$-gulosa}
    \begin{pseudo}
      \pl{Inicialize $Q(s,a)$ para todo $s \in \gS,\, a \in \gA$; $t=0$; estado inicial $s_0$}
      \pl{Defina $\polb \leftarrow \epsilon\text{-guloso}(Q)$ \hfill \textit{// política de comportamento}}
      \pl{\textbf{repita}}
      \pl[1.2em]{tome $a_t \sim \polb(s_t)$ e observe $(r_t, s_{t+1})$}
      \pl[1.2em]{$Q(s_t,a_t) \mathrel{+}= \alpha\bigl(r_t + \gamma \max_{a'} Q(s_{t+1},a') - Q(s_t,a_t)\bigr)$}
      \pl[1.2em]{$\polb \leftarrow \epsilon\text{-guloso}(Q)$; \quad $t \mathrel{+}= 1$}
      \pl{\textbf{fim}}
    \end{pseudo}
  \end{algbox}

\end{frame}

%% ------------------------------------------------------------
\begin{frame}{Convergência do Q-learning tabular}

  \begin{teobox}{Teorema}
    Para MDPs com $\gS$ e $\gA$ finitos, o Q-learning converge para o valor
    de ação ótimo, $Q(s,a) \to \Qstar(s,a)$, desde que \textbf{(i)} a
    sequência de políticas satisfaça a condição GLIE e \textbf{(ii)} os
    passos satisfaçam as condições de Robbins--Monro:
    \vspace{-0.7em}
    \[ \sum_{t=1}^{\infty} \alpha_t = \infty, \qquad
       \sum_{t=1}^{\infty} \alpha_t^2 < \infty. \]
  \end{teobox}

  \begin{ideiabox}
    A primeira soma garante que qualquer erro inicial \emph{pode} ser
    corrigido (os passos não somem cedo demais); a segunda garante que o
    ruído das amostras \emph{é} amortecido. Ex.: $\alpha_t = 1/t$ serve.
  \end{ideiabox}

  {\small A prova apoia-se em aproximação estocástica, na contração do
  backup de Bellman e em valores limitados. \alert{Nenhuma das três
  sobrevive intacta à aproximação de função.}}
\end{frame}

\begin{frame}{SARSA {\normalfont vs.} Q-learning no mesmo passo}
  \framesubtitle{\emph{Mars Rover} com $\alpha = 0{,}5$, $\gamma = 1$, $Q(s_7,a_1) = 10$ e $Q(s_7,a_2) = 5$}

  \begin{columns}[T]
    \begin{column}{0.5\textwidth}
      \textbf{SARSA} --- tupla $(s_6, a_1, 0, s_7, a_2, \dots)$:
      \vspace{-0.4em}
      \[ \small Q(s_6,a_1) \leftarrow 0 + \alpha\bigl(0 + \gamma\, Q(s_7,\mathbf{a_2})\bigr) = \mathbf{2{,}5} \]
      \vspace{-1.4em}

      {\footnotesize Usa $a_2$: foi a ação de fato tomada em $s_7$,
      possivelmente por exploração.}
    \end{column}
    \begin{column}{0.5\textwidth}
      \textbf{Q-learning} --- tupla $(s_6, a_1, 0, s_7)$:
      \vspace{-0.4em}
      \[ \small Q(s_6,a_1) \leftarrow 0 + \alpha\bigl(0 + \gamma \max_{a'} Q(s_7,a')\bigr) = \mathbf{5} \]
      \vspace{-1.4em}

      {\footnotesize Usa $\max\{10,5\} = 10$, ignorando a ação que o agente
      de fato escolheu.}
    \end{column}
  \end{columns}

  \begin{alertabox}
    A inicialização de $Q$ importa? \textbf{Assintoticamente não}, sob as
    condições do teorema anterior --- mas \textbf{muito} no início.
    Inicializar otimisticamente é, aliás, uma estratégia de exploração.
  \end{alertabox}

  \begin{ideiabox}
    A diferença é de um símbolo, mas é ela que separa ``qual o valor de
    seguir esta política, com exploração e tudo'' de ``qual o valor de
    jogar perfeitamente a partir daqui''.
  \end{ideiabox}
\end{frame}

\begin{frame}{O que essa diferença provoca: o \emph{cliff walking}}
  \framesubtitle{Sutton \& Barto (2018), Exemplo 6.6 --- o exemplo canônico da distinção}

  \begin{columns}[T]
    \begin{column}{0.56\textwidth}
      \begin{center}
      \scalebox{0.95}{\begin{tikzpicture}[x=7mm, y=7mm, font=\scriptsize]
        \foreach \x in {0,...,9}{\foreach \y in {0,...,2}{
          \draw[rlAzul!30, fill=rlAzul!5] (\x,\y) rectangle ++(1,1);}}
        \foreach \x in {1,...,8}{
          \draw[rlVermelho!55, fill=rlVermelho!18] (\x,0) rectangle ++(1,1);}
        \draw[rlVerde!60!black, fill=rlVerde!18] (0,0) rectangle ++(1,1);
        \draw[rlVerde!60!black, fill=rlVerde!18] (9,0) rectangle ++(1,1);
        \node at (0.5,0.5) {I};
        \node at (9.5,0.5) {M};
        \node[text=rlVermelho!85!black] at (4.5,0.5) {penhasco};
        \draw[rlLaranja, very thick, -{Stealth[length=2mm]}, rounded corners=2pt]
          (0.5,0.8) -- (0.5,1.5) -- (9.5,1.5) -- (9.5,0.82);
        \draw[rlAzul, very thick, -{Stealth[length=2mm]}, rounded corners=2pt]
          (0.32,0.8) -- (0.32,2.5) -- (9.68,2.5) -- (9.68,0.82);
        \node[text=rlLaranja, fill=white, inner sep=1pt, anchor=west] at (2.4,1.72) {Q-learning: rota ótima};
        \node[text=rlAzul, fill=white, inner sep=1pt, anchor=west] at (2.4,2.72) {SARSA: rota segura};
      \end{tikzpicture}}
      \end{center}
      {\footnotesize $r = -1$ por passo; cair no penhasco custa $-100$ e
      devolve o agente a I. Se $\epsilon \to 0$, ambos convergem para a
      política ótima --- mas \emph{durante} o aprendizado o comportamento
      cauteloso do SARSA costuma ser preferível.}
    \end{column}
    \begin{column}{0.44\textwidth}
      \begin{itemize}\setlength{\itemsep}{0.3ex}
        \item {\color{rlLaranja}\textbf{Q-learning}} aprende o caminho
          \emph{ótimo}, rente ao penhasco --- mas, agindo com
          $\epsilon$-guloso, cai nele de vez em quando e obtém
          \textbf{recompensa média pior}.
        \item {\color{rlAzul}\textbf{SARSA}} aprende o caminho \emph{seguro},
          pela rota de cima, porque seu alvo já contabiliza os passos
          exploratórios.
      \end{itemize}

    \end{column}
  \end{columns}
\end{frame}

%% ------------------------------------------------------------
\begin{frame}{Teste sua compreensão: SARSA e Q-learning}

  \begin{quizbox}{}
    Quais afirmações são verdadeiras?
    \begin{enumerate}\setlength{\itemsep}{0.2ex}
      \item Tanto SARSA quanto Q-learning podem atualizar a política após
        cada passo.
      \item Se $\epsilon = 0$ em todos os instantes e $Q$ é inicializado
        aleatoriamente, a atualização de SARSA coincide com a de Q-learning.
    \end{enumerate}
  \end{quizbox}
  \paradiscussao

  \pause
  \begin{respbox}
    \textbf{As duas são verdadeiras.} (1) Ambos são incrementais e a
    política gulosa muda assim que $Q$ muda. (2) Com $\epsilon = 0$ a
    política é gulosa pura, logo $a_{t+1} = \argmax_{a'} Q(s_{t+1},a')$ e
    os dois alvos coincidem. O que se perde é a garantia de convergência,
    que exige GLIE.
  \end{respbox}
\end{frame}

%% ------------------------------------------------------------
\begin{frame}{Um viés escondido no $\max$}
  \framesubtitle{Expansão: por que o Q-learning superestima valores}

  \begin{itemize}\setlength{\itemsep}{0.35ex}
    \item O alvo do Q-learning usa $\max_{a'} \Qhat(s',a')$ como estimativa
      de $\max_{a'} Q(s',a')$. Mas o máximo de estimativas ruidosas é um
      estimador \textbf{enviesado para cima}:
      \[ \E\Bigl[\max_{a'} \Qhat(s',a')\Bigr] \;\ge\; \max_{a'} \E\bigl[\Qhat(s',a')\bigr] \]
      pela desigualdade de Jensen ($\max$ é convexo).
    \item Basta que uma ação tenha tido sorte nas amostras para que o
      $\max$ a selecione --- e o erro é propagado adiante pelo
      \emph{bootstrapping}.
  \end{itemize}


  \begin{ideiabox}
    A correção clássica é \textbf{desacoplar seleção e avaliação} da ação:
    usar um conjunto de estimativas para escolher $\argmax$ e outro para
    avaliá-lo. É exatamente a ideia do \emph{Double Q-learning} --- e,
    depois, do \emph{Double DQN}, no fim da aula.
  \end{ideiabox}
\end{frame}

%% ============================================================
\section{Aproximação de função valor}
%% ============================================================

\begin{frame}{Por que a tabela não basta}
  \framesubtitle{O momento em que $\abs{\gS}$ deixa de ser um número que se escreve}

  \begin{columns}[T]
    \begin{column}{0.48\textwidth}
      \begin{center}\small
      \begin{tabular}{lr}
        \toprule
        Domínio & $\abs{\gS}$ aproximado \\
        \midrule
        \emph{Mars Rover} (Aula 2) & $7$ \\
        Gamão & $10^{20}$ \\
        Go ($19\times19$) & $10^{170}$ \\
        Atari (4 quadros $84{\times}84$) & $256^{28224}$ \\
        Robô contínuo & incontável \\
        \bottomrule
      \end{tabular}
      \end{center}
    \end{column}
    \begin{column}{0.52\textwidth}
      \begin{alertabox}
        Não é só memória: mesmo que a tabela coubesse, seria preciso
        \emph{visitar} cada entrada ao menos uma vez para estimá-la.
      \end{alertabox}

      \begin{ideiabox}
        A palavra-chave é \textbf{generalização}: aprender sobre um estado
        e, com isso, saber algo sobre estados parecidos que nunca foram
        visitados.
      \end{ideiabox}
    \end{column}
  \end{columns}

  \smallskip
  {\small Queremos evitar armazenar por estado o modelo $(P,\rec)$, o valor
  $V$, o valor de ação $Q$ e a própria política $\pol$ --- reduzindo
  \textbf{memória}, \textbf{computação} e, sobretudo, a
  \textbf{experiência} necessária.}
\end{frame}

\begin{frame}{O problema, primeiro com um oráculo}
  \framesubtitle{Reduzir RL a aprendizado supervisionado antes de complicar}

  \begin{itemize}\setlength{\itemsep}{0.35ex}
    \item Suponha, por ora, que exista um oráculo capaz de devolver o valor
      verdadeiro $Q^{\pol}(s,a)$ para qualquer par consultado.
    \item Então temos um problema de regressão comum, sobre pares
      $\bigl((s,a),\, Q^{\pol}(s,a)\bigr)$.
    \item Objetivo: achar $\wv$ na família $\Qhat(s,a;\wv)$ que melhor
      aproxima $Q^{\pol}$.
  \end{itemize}

  \vspace{-0.7em}
  \begin{defbox}{Função de perda}
    \[ J(\wv) \;=\; \E_{\pol}\Bigl[\bigl(Q^{\pol}(s,a) - \Qhat(s,a;\wv)\bigr)^2\Bigr] \]
  \end{defbox}

  \vspace{-0.7em}
  {\small Como
  $\nabla_{\wv} J(\wv) = -2\,\E_{\pol}\bigl[(Q^{\pol}(s,a) - \Qhat(s,a;\wv))\,\nabla_{\wv}\Qhat(s,a;\wv)\bigr]$,
  o passo $\Delta \wv = -\tfrac{1}{2}\alpha \nabla_{\wv} J(\wv)$ por amostra fica}
  \vspace{-0.6em}
  \[ \Delta \wv = \alpha\bigl(Q^{\pol}(s,a) - \Qhat(s,a;\wv)\bigr)\nabla_{\wv}\Qhat(s,a;\wv). \]
  \vspace{-1.3em}

  {\small \textbf{SGD}: usar uma (ou poucas) amostras no lugar da
  esperança. Em média, o passo é o mesmo do gradiente completo.}
\end{frame}

%% ------------------------------------------------------------
\begin{frame}{Aproximadores lineares e o caso tabular}
  \framesubtitle{Expansão: por que o tabular é um caso particular, e não uma alternativa}

  \begin{itemize}\setlength{\itemsep}{0.35ex}
    \item Descreva cada par por um vetor de atributos
      $\xv(s,a) = [x_1(s,a), \dots, x_d(s,a)]^{\top}$ e tome
      \[ \Qhat(s,a;\wv) = \wv^{\top}\xv(s,a), \qquad
         \nabla_{\wv}\Qhat(s,a;\wv) = \xv(s,a). \]
    \item A atualização vira, de forma particularmente limpa,
      \[ \Delta\wv = \alpha\bigl(\text{alvo} - \wv^{\top}\xv(s,a)\bigr)\,\xv(s,a). \]
  \end{itemize}

  \vspace{-0.3em}
  \begin{defbox}{O tabular como caso especial}
    Tome $d = \abs{\gS}\abs{\gA}$ e atributos \emph{one-hot}:
    $x_i(s,a) = \mathbf{1}\{(s,a) = i\}$. Então $\wv$ \emph{é} a tabela,
    $\nabla_{\wv}\Qhat$ seleciona uma única coordenada, e a atualização
    acima reduz-se exatamente ao Q-learning tabular. Não há dois
    algoritmos: há um só, com representações diferentes --- no caso
    one-hot simplesmente não existe generalização.
  \end{defbox}

\end{frame}

%% ------------------------------------------------------------
\begin{frame}{Sem oráculo: Monte Carlo com aproximação de função}

  {\small O retorno $G_t$ é amostra \textbf{não-enviesada}, ainda que
  ruidosa, de $Q^{\pol}(s_t,a_t)$: basta substituir o oráculo por $G_t$ e
  fazer regressão sobre os pares $\langle (s_t,a_t), G_t\rangle$. Com
  aproximação linear, isso converge ao mínimo global da perda quadrática.}

  \begin{algbox}{Avaliação de política MC com aproximação de função}
    \begin{pseudo}
      \pl{Inicialize $\wv$; $k = 1$}
      \pl{\textbf{repita}}
      \pl[1.2em]{gere o episódio $k$ seguindo $\pol$}
      \pl[1.2em]{\textbf{para} $t = 1,\dots,L_k$ \textbf{faça}}
      \pl[2.4em]{\textbf{se} primeira visita a $(s_t,a_t)$ \textbf{então}}
      \pl[3.6em]{$G_t = \sum_{j=t}^{L_k} \gamma^{\,j-t} r_{k,j}$;\;\;
                 $\wv \mathrel{+}= \alpha\bigl(G_t - \Qhat(s_t,a_t;\wv)\bigr)\nabla_{\wv}\Qhat(s_t,a_t;\wv)$}
      \pl[2.4em]{\textbf{fim}}
      \pl[1.2em]{\textbf{fim}; \quad $k \mathrel{+}= 1$}
      \pl{\textbf{fim}}
    \end{pseudo}
  \end{algbox}

\end{frame}

%% ------------------------------------------------------------
\begin{frame}{TD(0) com aproximação de função}
  \framesubtitle{Três aproximações empilhadas, cada uma com seu preço}

  \begin{itemize}\setlength{\itemsep}{0.3ex}
    \item No caso tabular, TD(0) atualizava
      $\Vpi(s) \leftarrow \Vpi(s) + \alpha(r + \gamma \Vpi(s') - \Vpi(s))$,
      com o alvo $r + \gamma\Vpi(s')$ já \emph{enviesado}.
    \item Com aproximação, o alvo passa a ser
      $r + \gamma\Vhat^{\pol}(s';\wv)$: enviesado \textbf{e} aproximado.
  \end{itemize}

  \begin{alertabox}
    \begin{enumerate}\setlength{\itemsep}{0.2ex}
      \item \textbf{Amostragem} --- usamos uma transição no lugar da
        esperança sobre $P$.
      \item \textbf{Bootstrapping} --- usamos nossa própria estimativa no
        lugar do valor verdadeiro.
      \item \textbf{Aproximação de função} --- projetamos o resultado numa
        família paramétrica que talvez não contenha $\Vpi$.
    \end{enumerate}
  \end{alertabox}

  \vspace{-0.3em}
  \begin{itemize}\setlength{\itemsep}{0.3ex}
    \item A perda otimizada é
      $J(\wv) = \E_{\pol}\bigl[(r_j + \gamma\Vhat(s_{j+1};\wv) - \Vhat(s_j;\wv))^2\bigr]$.
    \item \textbf{Detalhe importante:} tratamos o alvo como \emph{constante}
      ao derivar. Por isso a atualização não é o gradiente verdadeiro de
      $J$ --- chama-se \textbf{semigradiente}.
  \end{itemize}
\end{frame}

%% ------------------------------------------------------------
\begin{frame}{TD(0) com aproximação: o algoritmo}

  \begin{algbox}{Avaliação de política TD(0) com aproximação de função}
    \begin{pseudo}
      \pl{Inicialize $\wv$ e o estado $s$}
      \pl{\textbf{repita}}
      \pl[1.2em]{amostre $a \sim \pol(s)$; observe $r(s,a)$ e $s' \sim P(\cdot\mid s,a)$}
      \pl[1.2em]{$\Delta\wv = \alpha\bigl(r + \gamma\Vhat(s';\wv) - \Vhat(s;\wv)\bigr)\nabla_{\wv}\Vhat(s;\wv)$;\;\; $\wv \mathrel{+}= \Delta\wv$}
      \pl[1.2em]{\textbf{se} $s'$ não é terminal \textbf{então} $s \leftarrow s'$}
      \pl[1.2em]{\textbf{senão} reinicie o episódio, amostrando novo estado inicial $s$}
      \pl{\textbf{fim}}
    \end{pseudo}
  \end{algbox}

  \begin{ideiabox}
    Derivar também o alvo produziria o \emph{gradiente residual}: estável,
    porém lento e exigindo duas amostras independentes do próximo estado.
    O semigradiente aprende mais rápido --- ao custo das garantias.
  \end{ideiabox}


\end{frame}

%% ============================================================
\section{Controle com aproximação e Deep Q-Learning}
%% ============================================================

\begin{frame}{Controle com aproximação de função valor}

  \begin{itemize}\setlength{\itemsep}{0.35ex}
    \item Representamos os valores de ação por $\Qhat^{\pol}(s,a;\wv) \approx Q^{\pol}$.
    \item Intercalamos, como sempre:
      \begin{itemize}\setlength{\itemsep}{0.15ex}
        \item avaliação aproximada, ajustando $\wv$;
        \item melhoria $\epsilon$-gulosa sobre $\Qhat$.
      \end{itemize}
    \item Como não conhecemos $Q^{\pol}(s_t,a_t)$, substituímos por um
      \textbf{alvo}, exatamente como no caso tabular:
  \end{itemize}

  \smallskip
  {\small
  \begin{tabular}{@{}ll@{}}
    \textbf{Molde geral} & $\Delta\wv = \alpha\bigl(\text{alvo} - \Qhat(s_t,a_t;\wv)\bigr)\nabla_{\wv}\Qhat(s_t,a_t;\wv)$\\[0.7ex]
    \textbf{Monte Carlo} & alvo $= G_t$ \\[0.5ex]
    \textbf{SARSA} & alvo $= r + \gamma\,\Qhat(s',a';\wv)$ \\[0.5ex]
    \textbf{Q-learning} & alvo $= r + \gamma \max_{a'} \Qhat(s',a';\wv)$ \\
  \end{tabular}}

  \smallskip
  {\small Uma única linha de código, três algoritmos: toda a Seção 3
  reaparece trocando ``entrada da tabela'' por ``saída da rede''.}
\end{frame}

%% ------------------------------------------------------------
\begin{frame}{A ``tríade mortal''}
  \framesubtitle{Quando os três se encontram, a convergência deixa de ser garantida}

  \begin{columns}[T]
    \begin{column}{0.42\textwidth}
      \begin{center}
      \scalebox{0.86}{\begin{tikzpicture}[font=\scriptsize]
        \fill[rlVermelho, fill opacity=0.15] (0,0.7) circle (1.2);
        \fill[rlLaranja,  fill opacity=0.15] (-0.85,-0.75) circle (1.2);
        \fill[rlAmbar,    fill opacity=0.22] (0.85,-0.75) circle (1.2);
        \draw[rlVermelho, thick] (0,0.7) circle (1.2);
        \draw[rlLaranja, thick] (-0.85,-0.75) circle (1.2);
        \draw[rlAmbar!85!black, thick] (0.85,-0.75) circle (1.2);
        \node[align=center, text=rlVermelho!75!black, font=\scriptsize\bfseries]
          at (0,1.3) {aproximação\\de função};
        \node[align=center, text=rlLaranja!85!black, font=\scriptsize\bfseries]
          at (-1.35,-1.15) {boot-\\strapping};
        \node[align=center, text=rlAmbar!45!black, font=\scriptsize\bfseries]
          at (1.4,-1.15) {off-policy};
        \node[font=\scriptsize\bfseries, text=rlVermelho!80!black] at (0,-0.32) {instável};
      \end{tikzpicture}}
      \end{center}
    \end{column}
    \begin{column}{0.58\textwidth}
      \begin{itemize}\setlength{\itemsep}{0.25ex}
        \item Cada atualização faz um backup de Bellman aproximado e
          \emph{projeta} o resultado sobre a família de funções.
        \item O backup é uma \textbf{contração}; a projeção pode ser uma
          \textbf{expansão}. A composição pode não contrair --- e aí a
          sequência oscila ou diverge.
        \item Dois a dois são seguros; os três juntos, não.
      \end{itemize}
      \vspace{-0.5em}

      \begin{alertabox}
        Contraexemplo canônico: o de \textbf{Baird} (Sutton \& Barto,
        2018), em que Q-learning linear diverge com os pesos crescendo sem
        limite.
      \end{alertabox}
    \end{column}
  \end{columns}
\end{frame}

%% ------------------------------------------------------------
\begin{frame}{Mapa das garantias de convergência}
  \framesubtitle{Expansão: o que se sabe, para avaliação e para controle}

  \begin{center}\small
  \begin{tabular}{llccc}
    \toprule
    & Algoritmo & Tabular & Linear & Não-linear \\
    \midrule
    \multirow{3}{*}{\textbf{Avaliação}}
      & MC              & converge & converge & converge \\
      & TD(0) on-policy & converge & converge$^{\dagger}$ & pode divergir \\
      & TD(0) off-policy& converge & pode divergir & pode divergir \\
    \midrule
    \multirow{3}{*}{\textbf{Controle}}
      & MC-controle     & converge & oscila próximo & pode divergir \\
      & SARSA           & converge & oscila próximo & pode divergir \\
      & Q-learning      & converge & pode divergir & pode divergir \\
    \bottomrule
  \end{tabular}
  \end{center}

  {\footnotesize $^{\dagger}$ para um ponto fixo cujo erro é limitado por
  $\frac{1}{1-\gamma}$ vezes o erro da melhor aproximação da família.}

  \vspace{-0.5em}
  \begin{ideiabox}
    Note onde está o Q-learning com rede neural: sem garantia alguma. O DQN
    não resolve isso teoricamente --- ele \emph{engenheira} o problema para
    que, na prática, funcione.
  \end{ideiabox}
\end{frame}

%% ------------------------------------------------------------
\begin{frame}{Q-learning com redes neurais: o que dá errado}

  {\small No caso tabular o Q-learning converge para $\Qstar$; com
  aproximação, ele minimiza um EQM por SGD usando um alvo estimado --- e
  pode divergir.}

  \begin{alertabox}
    \textbf{Dois problemas concretos, além da tríade:}
    \begin{enumerate}\setlength{\itemsep}{0.25ex}
      \item \textbf{Correlação entre amostras.} Transições consecutivas são
        dependentes; SGD supõe amostras aproximadamente i.i.d.
      \item \textbf{Alvos não-estacionários.} O alvo usa os mesmos pesos
        que estamos atualizando: perseguimos um alvo que se move.
    \end{enumerate}
  \end{alertabox}

  \begin{defbox}{As duas respostas do DQN \textnormal{(Mnih et al., 2015)}}
    \textbf{Repetição de experiência} (\emph{experience replay}) contra a
    correlação; \textbf{alvos fixos} (\emph{fixed Q-targets}) contra a
    não-estacionariedade.
  \end{defbox}
\end{frame}

%% ------------------------------------------------------------
\begin{frame}{Repetição de experiência}
  \framesubtitle{Guardar o passado e reamostrá-lo fora de ordem}

  \begin{itemize}\setlength{\itemsep}{0.3ex}
    \item Mantenha um \textbf{buffer de repetição} $\mathcal{D}$ com as
      tuplas $(s,a,r,s')$ já vividas (tipicamente $10^6$ transições, em
      janela deslizante).
    \item Repita:
      \begin{enumerate}\setlength{\itemsep}{0.15ex}
        \item amostre $(s,a,r,s') \sim \mathcal{D}$ (ou um \emph{minibatch});
        \item calcule o alvo $r + \gamma \max_{a'} \Qhat(s',a';\wv)$;
        \item atualize
          $\Delta\wv = \alpha\bigl(r + \gamma\max_{a'}\Qhat(s',a';\wv) - \Qhat(s,a;\wv)\bigr)\nabla_{\wv}\Qhat(s,a;\wv)$.
      \end{enumerate}
  \end{itemize}

  \begin{ideiabox}
    Só é possível porque o Q-learning é \textbf{off-policy}: as tuplas do
    buffer foram geradas por políticas antigas, e isso não invalida o alvo
    (que já usa o $\max$, não a ação realmente escolhida). SARSA não pode
    fazer o mesmo sem correção. Resta, porém, o segundo problema: o alvo é
    tratado como escalar, mas muda quando os pesos mudam.
  \end{ideiabox}

\end{frame}

%% ------------------------------------------------------------
\begin{frame}{Alvos fixos}
  \framesubtitle{Congelar a rede que produz o alvo por $C$ atualizações}

  \begin{itemize}\setlength{\itemsep}{0.35ex}
    \item Mantenha \textbf{dois} conjuntos de pesos:
      \begin{itemize}\setlength{\itemsep}{0.15ex}
        \item $\wv$ --- os pesos que estão sendo atualizados;
        \item $\wtar$ --- os pesos da \emph{rede-alvo}, uma cópia congelada.
      \end{itemize}
    \item O alvo passa a ser calculado com $\wtar$:
      \[ \Delta\wv = \alpha\Bigl(r + \gamma \max_{a'} \Qhat(s',a';\mdestaque{rlLaranja}{\wtar}) - \Qhat(s,a;\wv)\Bigr)\nabla_{\wv}\Qhat(s,a;\wv) \]
    \item A cada $C$ passos, sincronize: $\wtar \leftarrow \wv$.
  \end{itemize}

  \begin{ideiabox}
    Entre duas sincronizações, o problema vira uma regressão supervisionada
    comum, com rótulos fixos. É a manobra que devolve estabilidade a um
    procedimento que, de outro modo, persegue a própria cauda.
  \end{ideiabox}
\end{frame}

%% ------------------------------------------------------------
\begin{frame}{DQN: o algoritmo completo}

  \begin{algbox}{Deep Q-Network \textnormal{(Mnih et al., 2015)}}
    \begin{pseudo}
      \pl{Entrada: $C$, $\alpha$, buffer $\mathcal{D} = \{\}$; inicialize $\wv$, $\wtar \leftarrow \wv$, $t = 0$; obtenha $s_0$}
      \pl{\textbf{repita}}
      \pl[1.2em]{amostre $a_t$ pela política $\epsilon$-gulosa sobre $\Qhat(s_t,\cdot;\wv)$}
      \pl[1.2em]{observe $r_t$ e $s_{t+1}$; armazene $(s_t,a_t,r_t,s_{t+1})$ em $\mathcal{D}$}
      \pl[1.2em]{amostre um \emph{minibatch} aleatório $\{(s_i,a_i,r_i,s_{i+1})\}$ de $\mathcal{D}$}
      \pl[1.2em]{\textbf{para} cada $i$ no \emph{minibatch} \textbf{faça}}
      \pl[2.4em]{$y_i = r_i$ \; se o episódio terminou em $i+1$;\quad
                 $y_i = r_i + \gamma \max_{a'} \Qhat(s_{i+1},a';\wtar)$ \; caso contrário}
      \pl[2.4em]{passo de descida em $\bigl(y_i - \Qhat(s_i,a_i;\wv)\bigr)^2$:\;\;
                 $\Delta\wv = \alpha\bigl(y_i - \Qhat(s_i,a_i;\wv)\bigr)\nabla_{\wv}\Qhat(s_i,a_i;\wv)$}
      \pl[1.2em]{\textbf{fim};\quad $t \mathrel{+}= 1$}
      \pl[1.2em]{\textbf{se} $t \bmod C = 0$ \textbf{então} $\wtar \leftarrow \wv$}
      \pl{\textbf{fim}}
    \end{pseudo}
  \end{algbox}

  {\footnotesize Há várias escolhas de projeto embutidas: arquitetura da
  rede, taxa de aprendizado, frequência $C$ de sincronização, tamanho do
  buffer e como preenchê-lo antes de começar a aprender.}
\end{frame}

%% ------------------------------------------------------------
\begin{frame}{Teste sua compreensão: alvos fixos}

  \begin{quizbox}{}
    No DQN o alvo usa um conjunto separado de pesos,
    $r + \gamma\max_{a'}\Qhat(s',a';\wtar)$. Quais afirmações são
    verdadeiras?
    \begin{enumerate}\setlength{\itemsep}{0.2ex}
      \item Isso \emph{dobra o tempo de computação} em relação a um método
        sem rede-alvo separada.
      \item Isso \emph{dobra o requisito de memória} em relação a um método
        sem rede-alvo separada.
    \end{enumerate}
  \end{quizbox}
  \paradiscussao

  \pause
  \begin{respbox}
    \textbf{Apenas a segunda.} A memória dobra: são dois conjuntos de pesos
    armazenados. O tempo \emph{não} dobra --- em ambos os casos há uma
    passagem para frente (alvo) e uma para frente/para trás (atualização).
    O extra é só a cópia $\wtar \leftarrow \wv$ a cada $C$ passos.
  \end{respbox}
\end{frame}

%% ------------------------------------------------------------
\begin{frame}{DQN em Atari}
  \framesubtitle{Aprendizado ponta a ponta, de pixels a valores de ação}

  \begin{columns}[T]
    \begin{column}{0.55\textwidth}
      \begin{center}
      \scalebox{0.95}{\begin{tikzpicture}[node distance=6mm, font=\scriptsize]
        \node[draw=rlAzul, thick, fill=rlAzul!8, minimum height=13mm,
              minimum width=15mm, align=center, rounded corners=1pt] (px)
              {4 quadros\\$84\times84$};
        \node[draw=rlAzulClaro, thick, fill=rlAzulClaro!10, minimum height=13mm,
              minimum width=15mm, align=center, right=of px, rounded corners=1pt] (cv)
              {3 camadas\\convolucionais};
        \node[draw=rlTeal, thick, fill=rlTeal!10, minimum height=13mm,
              minimum width=13mm, align=center, right=of cv, rounded corners=1pt] (fc)
              {camada\\densa};
        \node[draw=rlLaranja, very thick, fill=rlLaranja!12, minimum height=13mm,
              minimum width=17mm, align=center, right=of fc, rounded corners=1pt] (out)
              {$\Qhat(s,a_1)\;\dots$\\$\Qhat(s,a_{18})$};
        \foreach \a/\b in {px/cv, cv/fc, fc/out}{\draw[trans] (\a) -- (\b);}
      \end{tikzpicture}}
      \end{center}
    \end{column}
    \begin{column}{0.45\textwidth}
      \begin{itemize}\setlength{\itemsep}{0.25ex}
        \item Estado $s$: pilha dos \textbf{4 últimos quadros} de pixels
          brutos --- para recuperar velocidade e direção, que um quadro
          único não revela.
        \item Saída: um valor $\Qhat(s,a)$ para cada uma das 18 posições de
          joystick, em uma única passagem.
        \item Recompensa: variação de pontuação (recortada em $[-1,1]$).
        \item \textbf{A mesma arquitetura e os mesmos hiperparâmetros} em
          todos os jogos.
      \end{itemize}
    \end{column}
  \end{columns}

\end{frame}

%% ------------------------------------------------------------
\begin{frame}{Quais ingredientes do DQN realmente importavam?}
  \framesubtitle{Estudo de ablação de Mnih et al.\ (2015): pontuação por jogo}

  \begin{center}\small
  \begin{tabular}{lrrrrr}
    \toprule
    Jogo & Linear & Rede & DQN & DQN & DQN \\
         &        & profunda & alvo fixo & \emph{replay} & \emph{replay} + alvo \\
    \midrule
    Breakout       & 3    & 3    & 10   & 241  & \textbf{317} \\
    Enduro         & 62   & 29   & 141  & 831  & \textbf{1006} \\
    River Raid     & 2345 & 1453 & 2868 & 4102 & \textbf{7447} \\
    Seaquest       & 656  & 275  & 1003 & 823  & \textbf{2894} \\
    Space Invaders & 301  & 302  & 373  & 826  & \textbf{1089} \\
    \bottomrule
  \end{tabular}
  \end{center}

  \begin{alertabox}
    Trocar o modelo linear por uma rede profunda, \textbf{sozinho}, chegou
    a \emph{piorar} o desempenho (Enduro: 62 $\to$ 29; Seaquest: 656 $\to$
    275). O que fez a diferença foram o \emph{replay} e o alvo fixo.
  \end{alertabox}

\end{frame}

\begin{frame}{Por que o \emph{replay} ajuda tanto}
  \framesubtitle{Três efeitos, não um}

  \begin{enumerate}\setlength{\itemsep}{0.5ex}
    \item \textbf{Descorrelaciona} as amostras dentro do \emph{minibatch},
      aproximando a hipótese i.i.d.\ que o SGD supõe.
    \item \textbf{Reaproveita} cada transição muitas vezes --- ganho enorme
      de eficiência de dados, que é o recurso caro quando cada amostra
      exige uma interação com o ambiente.
    \item \textbf{Estabiliza a distribuição de treino}: sem buffer, uma
      mudança de política altera imediatamente a distribuição dos dados,
      e a rede ``esquece'' o que aprendeu em regiões que deixou de visitar
      (esquecimento catastrófico).
  \end{enumerate}

  \begin{ideiabox}
    O terceiro efeito é o mais subestimado: o buffer age como uma memória
    que impede o agente de otimizar apenas para o presente. É por isso que
    o \emph{replay} sozinho já produz saltos de $3 \to 241$ em Breakout.
  \end{ideiabox}
\end{frame}

%% ------------------------------------------------------------
\begin{frame}{Aperfeiçoamentos imediatos do DQN}

  \begin{itemize}\setlength{\itemsep}{0.5ex}
    \item \textbf{Double DQN} (van Hasselt et al., AAAI 2016) --- corrige o
      viés de maximização: a rede em treinamento \emph{escolhe} a ação,
      a rede-alvo a \emph{avalia}:
      \[ y = r + \gamma\, \Qhat\bigl(s', \argmax_{a'} \Qhat(s',a';\wv);\, \wtar\bigr) \]
    \item \textbf{Repetição priorizada} (Schaul et al., ICLR 2016) ---
      amostrar do buffer com probabilidade proporcional ao erro TD, em vez
      de uniformemente: aprender mais onde ainda se erra mais.
    \item \textbf{Dueling DQN} (Wang et al., ICML 2016) --- decompor
      $\Qhat(s,a) = \Vhat(s) + \widehat{A}(s,a)$ em valor de estado e
      vantagem, aprendendo $\Vhat(s)$ mesmo quando a ação é irrelevante.
  \end{itemize}

  \begin{ideiabox}
    Os três são ortogonais e combináveis --- é o que faz o \emph{Rainbow}
    (Hessel et al., 2018), somando seis dessas extensões.
  \end{ideiabox}
\end{frame}

%% ------------------------------------------------------------
\begin{frame}{O que você deve levar desta aula}

  \begin{defbox}{Capacidades esperadas}
    \begin{itemize}\setlength{\itemsep}{0.3ex}
      \item Implementar avaliação de política por MC e TD(0), tabular e com
        aproximação de função.
      \item Implementar controle por MC, SARSA e Q-learning, e explicar a
        diferença entre os alvos de SARSA e de Q-learning.
      \item Enunciar as condições GLIE e Robbins--Monro e dizer para que
        cada uma serve.
      \item Listar os três elementos da tríade mortal --- aproximação de
        função, \emph{bootstrapping} e aprendizado off-policy --- e
        descrever qualitativamente o problema que geram.
      \item Explicar as duas ideias que tornaram o DQN viável: repetição de
        experiência e alvos fixos.
    \end{itemize}
  \end{defbox}

\end{frame}

%% ------------------------------------------------------------
\begin{frame}{Próxima aula}

  \begin{itemize}\setlength{\itemsep}{0.45ex}
    \item Até aqui, todo método aprendeu uma \textbf{função valor} e
      derivou a política dela por $\argmax$.
    \item Isso tem limites: ações contínuas tornam o $\argmax$ caro, e
      políticas estocásticas ótimas não são representáveis dessa forma.
    \item \textbf{Próxima aula: gradientes de política} --- parametrizar
      $\pol_\theta$ diretamente e subir o gradiente do retorno esperado.
  \end{itemize}

  \begin{center}
    \scalebox{0.9}{\begin{tikzpicture}[node distance=7mm, font=\scriptsize]
      \node[draw=rlAzul, thick, fill=rlAzul!8, rounded corners=2pt,
            minimum width=30mm, minimum height=8mm] (a) {baseado em valor (hoje)};
      \node[draw=rlCinza, thick, fill=rlFundo, rounded corners=2pt,
            minimum width=30mm, minimum height=8mm, right=14mm of a] (b) {ator-crítico};
      \node[draw=rlLaranja, very thick, fill=rlLaranja!12, rounded corners=2pt,
            minimum width=30mm, minimum height=8mm, right=14mm of b] (c) {gradiente de política};
      \draw[trans] (a) -- (b);
      \draw[trans] (b) -- (c);
    \end{tikzpicture}}
  \end{center}

  \begin{ideiabox}
    Guarde a tríade mortal: ela reaparece, com outra roupa, em quase todo
    algoritmo profundo do restante do curso.
  \end{ideiabox}
\end{frame}

\end{document}
